\begin{table}[p]
\centering
\begin{singlespace}
\small
\caption{Annotation-interval coverage and alarm burden of selected full-development systems. Zero matching tolerance.}\label{tab:engineering}
\begin{tabularx}{\linewidth}{L{32mm}YYYYY}
\toprule
Selector & \makecell[r]{Hits /\\intervals} & \makecell[r]{Non-event\\clusters} & \makecell[r]{Clusters per\\negative h} & \makecell[r]{Redundant\\clusters} & \makecell[r]{Negative\\exposure (h)} \\
\midrule
\multicolumn{6}{l}{\textit{PVS vehicle/driver}} \\
Random-window & 59/60 & 71 & 24.05 & 14 & 2.9522 \\
Matched & 59/60 & 71 & 24.05 & 14 & 2.9522 \\
\multicolumn{6}{l}{\textit{PVS scenario}} \\
Random-window & 34/60 & 22 & 7.45 & 14 & 2.9522 \\
Matched & 52/60 & 51 & 17.28 & 18 & 2.9522 \\
\midrule\multicolumn{6}{l}{\textit{Karlsruhe final GPS-date group}} \\
Random-window & 7/12 & 0 & 0.00 & 0 & 0.0579 \\
Buffered-date & 7/12 & 0 & 0.00 & 1 & 0.0579 \\
\bottomrule
\end{tabularx}
\par\smallskip
\raggedright Higher interval coverage and lower alarm rates are desirable. Negative exposure means explicit non-bump time in PVS and explicit good-labelled time in Karlsruhe. PVS whole-trip results equal random-window results. Each PVS axis reuses the same nine acquisitions; Karlsruhe has one final date and only 3.47 minutes of assessed good exposure. Intervals are labelled passages, with repeated physical features possible.
\end{singlespace}
\end{table}
